% ════════════════════════════════════════════
% 模型的评价（7.模型评价.tex）写作规则 —— 模板内嵌，AI 先读本文件再写
% ════════════════════════════════════════════
\section{模型的评价}

\subsection{模型的优点}
\begin{itemize}[itemindent=2em]
\item \textbf{质量评分有明确依据。} 表\ref{tab:p1_dir}逐项说明 22 个指标的方向，广告字段还通过原文抽查确认“第 1 类=无广告”。处理时将负向指标反向转换、长度指标转为中心接近度，8 个列表指标通过 softmax 期望压缩，并用中位数填补 19 处缺失。抽样与扩展集合的可比领域差值约为 $0.002$，为这两类数据上的评分一致性提供了对照。
\item \textbf{指标冲突可以相对比较。} 本文同时给出 $22.3\%$ 的冲突率、$0.069$ 的 Kendall 系数和 $28.9\%$ 的独立参照。在 $\tau\in[0.5,2.5]$ 的检验范围，比值均 $<0.92$。这样可以区分“存在共同倾向”与“高度一致”，再结合三类分歧原因和处理规则解释冲突样本。
\item \textbf{质量项便于解释和比较。} $C(1-Q)^{\gamma}$ 在 $Q=1$ 时归零，恢复 Chinchilla 的函数形式；与直接使用 $E-CQ^{\gamma}$ 相比，不会因减去质量项而在高质量端出现负损失。模型选择比较了四种形式的 $R^2$、残差信息评分及范围内外表现。B3 插值轨迹的数据指数约为 $-0.2805$，与 $-0.2799$ 接近，提供了内部一致性检查。
\item \textbf{配置变化可从成本公式解释。} 模型保留 $C_Q=D[g(Q)-g(Q_0)]_+$，使质量费用随处理的数据量变化，并使用题设附录 B 的三类函数参数。由此求得 $L_{\text{crit}}=6/\eta=30{,}000$，并推导式\eqref{eq:p3_balance}。折算量 $\tau(Q)$ 将质量费用换成与参数量相同的单位，便于理解质量投入如何改变参数与数据的平衡。
\item \textbf{不同来源的结果分别解释。} 问题四列明许可、时间变量和模型类型，在 C8 的 $1{,}860$ 个有效目录中，用主样本交集的 $964$ 个目录汇总 $37$ 项任务；六维指标另与完整 C1 作名称匹配和记录级秩相关检查。C4 用于宏观资源及字段覆盖分析，预测中区分重抽样区间和情景范围。对 B2/B8 的异常也保留检验说明，例如 B8 存在低于不可约损失下限的记录，避免把不同性质的数据当作同等证据。
\end{itemize}

\subsection{模型的缺点}
\begin{itemize}[itemindent=2em]
\item \textbf{配比关系仍是线性近似。} 模型没有领域交互项，13 个验证域的平均 $R^2$ 为 $0.629$，约 $37\%$ 的变异仍未解释，可能还与随机种子、数据顺序等因素有关。nih\_exporter、enron\_emails、europarl、philpapers 缺少自身损失列，优化时固定在参考占比，难以单独判断其影响。混合质量指数只带来 $9.9\times10^{-5}$ 的 $R^2$ 增量，说明当前线性组合没有提供明显的新信息，也限制了进一步利用 17 域评分细化配方的空间。
\item \textbf{部分评分方向仍依赖字段含义。} 广告、流畅度已用原文及领域对照检查，但 modernbert 的 6 级指标、qurater 的 4 级评级等，仍按“高等级对应高质量”的语义处理，缺少独立验证。熵权也受到分布形态影响，改用等权后领域排序一致性为 $0.679$。第\ref{sec:check}节还显示，方向误判会明显改变排名，因此排名细节应与方向和赋权方法一起解释。
\item \textbf{质量实验的支持范围有限。} 质量项主要依靠 B6/B7 半合成实验，B8 方向相反且含低于损失下限的记录，不能并入主拟合；B2 的 MAPE 为 $33.7\%$，也说明跨来源估计仍有较大偏差。质量配置只在 $[Q_0,1]$ 内讨论，更低质量的外推尚缺数据支持。
\item \textbf{能力预测受到时间和评分限制。} 损失映射的 $R^2=0.26$、MAPE 为 $67.8\%$，不能精确换算收益。C1 可比窗口约 10 个月，预测却长达其 $1.2\sim2.4$ 倍，C3 历史评测又存在来源差异。两个年份的回归不足以识别具体技术机制，非规模项还包含样本构成等因素。因此，长期规模占比 $50.5\%$ 与 2024 年 6—9 月的 $4.6\%$ 只能分别对应各自窗口。
\end{itemize}
